\documentclass[10pt]{extarticle}
\usepackage[a4paper,margin=1.5cm]{geometry}
\usepackage[T1]{fontenc}
\usepackage{amsmath,amssymb}
\usepackage{graphicx}
\usepackage[export]{adjustbox}
\usepackage{multicol}
\usepackage{enumitem}
\usepackage{booktabs,array,tabularx}
\setlength{\columnsep}{0.7cm}
\setlength{\parindent}{0pt}
\setlength{\parskip}{3pt}
\pagestyle{empty}
\makeatletter
% figure registry (figdefs.tex is generated by the extract step)
\newcommand{\qfigdef}[4]{\@namedef{qf@file@#1}{#2}\@namedef{qf@w@#1}{#3}\@namedef{qf@wide@#1}{#4}}
\newcommand{\qf@check}[1]{\@ifundefined{qf@file@#1}{\PackageError{score2}{Unknown figure id #1}{}}{}}
\newcommand{\qf@img}[4][]{\includegraphics[width=#2\dimexpr\@nameuse{qf@w@#4}pt\relax,max width=#3,max height=0.45\textheight,#1]{\@nameuse{qf@file@#4}}}
% Full-width questions: finalize wraps every question that contains a wide figure in \qfullbegin...\qfullend,
% so the WHOLE question (stem, figures, options) is set across the page in reading order, then the two columns
% resume. The question's head - its number, stem and first figure (\qfullsplit marks the end of the head) -
% is kept on one page, so a number is never left at a page foot with its figure overleaf; the rest may break
% across pages normally, so no big gaps are left. Without a \qfullsplit, the whole question is kept together
% if it is under 45% of a page. Never write these yourself.
\newif\ifqf@full
\newif\ifqf@boxed
\newsavebox{\qf@box}
\newcommand{\qf@open}{\global\qf@fulltrue\par\global\qf@boxedtrue
  \setbox\qf@box\vbox\bgroup\hsize\textwidth\linewidth\textwidth}
\newcommand{\qf@emit}[1]{\par\egroup\global\qf@boxedfalse
  \ifdim\dimexpr\ht\qf@box+\dp\qf@box\relax<#1\textheight\noindent\usebox{\qf@box}\par\else\unvbox\qf@box\fi}
\newcommand{\qfullsplit}{\ifqf@boxed\qf@emit{0.85}\fi}
\newcommand{\qf@close}{\ifqf@boxed\qf@emit{0.45}\fi\global\qf@fullfalse}
\newcommand{\qfullbegin}{\end{multicols}\qf@open}
\newcommand{\qfullend}{\qf@close\begin{multicols}{2}}
\newcommand{\qfullstart}{\qf@open}
\newcommand{\qfulldone}{\qf@close}
% \qfig[scale]{id}: block figure, centred, directly where it is written (size set by the tool).
% Fallback if a wide figure is not inside a full-width question: break out of the columns in place - no floats.
\newcommand{\qfig}[2][1]{\qf@check{#2}%
  \ifqf@full\qf@blk{#1}{#2}\else
  \ifnum\@nameuse{qf@wide@#2}=1 \expandafter\@firstoftwo\else\expandafter\@secondoftwo\fi
  {\end{multicols}\begin{center}\qf@img{#1}{\textwidth}{#2}\end{center}\begin{multicols}{2}}%
  {\qf@blk{#1}{#2}}\fi}
\newcommand{\qf@blk}[2]{\par\smallskip{\centering\qf@img{#1}{\linewidth}{#2}\par}\smallskip}
% \qfigi[scale]{id}: inline figure (inside option lists / tables / minipages); never breaks columns.
\newcommand{\qfigi}[2][1]{\qf@check{#2}\qf@img[valign=c]{#1}{\linewidth}{#2}}
\makeatother
% structure macros - lint reads these to check completeness and figure placement
\newcommand{\papersec}[1]{\par\bigskip{\centering\Large\bfseries PAPER-#1\par}\medskip}
\newcommand{\partsec}[2]{\par\medskip\noindent{\large\bfseries PART-#1: #2}\par\smallskip}
\newcommand{\sect}[1]{\par\smallskip\noindent\textbf{#1}\par}
\newcommand{\qn}[1]{\par\smallskip\noindent\textbf{#1.}~}
\newcommand{\paraq}[2]{\par\smallskip\noindent\textbf{Paragraph for Questions #1 and #2}\par}
\newenvironment{qoptions}{\begin{itemize}[leftmargin=1.8em,itemsep=1pt,topsep=2pt,label={}]}{\end{itemize}}
\newcommand{\opt}[2]{\item[(#1)] #2}
% figure registry (generated by score2_tools.py finalize)
\qfigdef{p1_phy_s1_q03_1}{Score Advanced - Test 16 (20 May 2023)-img/p1_phy_s1_q03_1.png}{73.5}{0}
\qfigdef{p1_phy_s1_q04_B}{Score Advanced - Test 16 (20 May 2023)-img/p1_phy_s1_q04_B.png}{177.0}{0}
\qfigdef{p1_phy_s2_q05_1}{Score Advanced - Test 16 (20 May 2023)-img/p1_phy_s2_q05_1.png}{119.0}{0}
\qfigdef{p1_phy_s2_q07_1}{Score Advanced - Test 16 (20 May 2023)-img/p1_phy_s2_q07_1.png}{137.3}{0}
\qfigdef{p1_phy_s2_q08_1}{Score Advanced - Test 16 (20 May 2023)-img/p1_phy_s2_q08_1.png}{82.0}{0}
\qfigdef{p1_phy_s2_q09_1}{Score Advanced - Test 16 (20 May 2023)-img/p1_phy_s2_q09_1.png}{68.6}{0}
\qfigdef{p1_phy_s2_q10_1}{Score Advanced - Test 16 (20 May 2023)-img/p1_phy_s2_q10_1.png}{154.8}{0}
\qfigdef{p1_phy_s3_para3_4_1}{Score Advanced - Test 16 (20 May 2023)-img/p1_phy_s3_para3_4_1.png}{103.0}{0}
\qfigdef{p1_phy_s4_q07_1}{Score Advanced - Test 16 (20 May 2023)-img/p1_phy_s4_q07_1.png}{204.0}{0}
\qfigdef{p1_phy_s4_q08_1}{Score Advanced - Test 16 (20 May 2023)-img/p1_phy_s4_q08_1.png}{97.7}{0}
\qfigdef{p1_phy_s4_q09_1}{Score Advanced - Test 16 (20 May 2023)-img/p1_phy_s4_q09_1.png}{135.1}{0}
\qfigdef{p1_chem_s1_q04_A}{Score Advanced - Test 16 (20 May 2023)-img/p1_chem_s1_q04_A.png}{52.3}{0}
\qfigdef{p1_chem_s1_q04_B}{Score Advanced - Test 16 (20 May 2023)-img/p1_chem_s1_q04_B.png}{51.7}{0}
\qfigdef{p1_chem_s1_q04_C}{Score Advanced - Test 16 (20 May 2023)-img/p1_chem_s1_q04_C.png}{51.7}{0}
\qfigdef{p1_chem_s1_q04_D}{Score Advanced - Test 16 (20 May 2023)-img/p1_chem_s1_q04_D.png}{52.4}{0}
\qfigdef{p1_chem_s2_q09_A}{Score Advanced - Test 16 (20 May 2023)-img/p1_chem_s2_q09_A.png}{112.7}{0}
\qfigdef{p1_chem_s2_q09_B}{Score Advanced - Test 16 (20 May 2023)-img/p1_chem_s2_q09_B.png}{168.5}{0}
\qfigdef{p1_chem_s2_q09_C}{Score Advanced - Test 16 (20 May 2023)-img/p1_chem_s2_q09_C.png}{168.7}{0}
\qfigdef{p1_chem_s2_q09_D}{Score Advanced - Test 16 (20 May 2023)-img/p1_chem_s2_q09_D.png}{168.9}{0}
\qfigdef{p1_chem_s3_q05_1}{Score Advanced - Test 16 (20 May 2023)-img/p1_chem_s3_q05_1.png}{86.7}{0}
\qfigdef{p1_chem_s3_q05_2}{Score Advanced - Test 16 (20 May 2023)-img/p1_chem_s3_q05_2.png}{47.2}{0}
\qfigdef{p1_chem_s4_q09_1}{Score Advanced - Test 16 (20 May 2023)-img/p1_chem_s4_q09_1.png}{289.9}{1}
\qfigdef{p2_phy_s1_q01_1}{Score Advanced - Test 16 (20 May 2023)-img/p2_phy_s1_q01_1.png}{108.0}{0}
\qfigdef{p2_phy_s1_q04_1}{Score Advanced - Test 16 (20 May 2023)-img/p2_phy_s1_q04_1.png}{71.1}{0}
\qfigdef{p2_phy_s1_q05_1}{Score Advanced - Test 16 (20 May 2023)-img/p2_phy_s1_q05_1.png}{164.6}{0}
\qfigdef{p2_phy_s1_q06_1}{Score Advanced - Test 16 (20 May 2023)-img/p2_phy_s1_q06_1.png}{118.2}{0}
\qfigdef{p2_phy_s2_para7_8_1}{Score Advanced - Test 16 (20 May 2023)-img/p2_phy_s2_para7_8_1.png}{125.6}{0}
\qfigdef{p2_phy_s2_para9_10_1}{Score Advanced - Test 16 (20 May 2023)-img/p2_phy_s2_para9_10_1.png}{138.5}{0}
\qfigdef{p2_phy_s3_para3_4_1}{Score Advanced - Test 16 (20 May 2023)-img/p2_phy_s3_para3_4_1.png}{327.5}{1}
\qfigdef{p2_phy_s3_para5_6_1}{Score Advanced - Test 16 (20 May 2023)-img/p2_phy_s3_para5_6_1.png}{245.2}{0}
\qfigdef{p2_phy_s4_q08_1}{Score Advanced - Test 16 (20 May 2023)-img/p2_phy_s4_q08_1.png}{111.1}{0}
\qfigdef{p2_phy_s4_q09_1}{Score Advanced - Test 16 (20 May 2023)-img/p2_phy_s4_q09_1.png}{87.5}{0}
\qfigdef{p2_chem_s1_q02_1}{Score Advanced - Test 16 (20 May 2023)-img/p2_chem_s1_q02_1.png}{73.2}{0}
\qfigdef{p2_chem_s2_para7_8_1}{Score Advanced - Test 16 (20 May 2023)-img/p2_chem_s2_para7_8_1.png}{181.8}{0}
\qfigdef{p2_chem_s2_q09_A}{Score Advanced - Test 16 (20 May 2023)-img/p2_chem_s2_q09_A.png}{83.5}{0}
\qfigdef{p2_chem_s2_q09_B}{Score Advanced - Test 16 (20 May 2023)-img/p2_chem_s2_q09_B.png}{83.2}{0}
\qfigdef{p2_chem_s2_q09_C}{Score Advanced - Test 16 (20 May 2023)-img/p2_chem_s2_q09_C.png}{83.8}{0}
\qfigdef{p2_chem_s2_q09_D}{Score Advanced - Test 16 (20 May 2023)-img/p2_chem_s2_q09_D.png}{83.2}{0}
\qfigdef{p2_chem_s3_q01_1}{Score Advanced - Test 16 (20 May 2023)-img/p2_chem_s3_q01_1.png}{71.3}{0}
\qfigdef{p2_chem_s3_q02_1}{Score Advanced - Test 16 (20 May 2023)-img/p2_chem_s3_q02_1.png}{134.4}{0}
\qfigdef{p2_chem_s4_q08_1}{Score Advanced - Test 16 (20 May 2023)-img/p2_chem_s4_q08_1.png}{185.2}{0}
\raggedcolumns
\makeatletter\@beginparpenalty=10000\makeatother
\newcolumntype{L}[1]{>{\raggedright\arraybackslash}p{#1}}
\begin{document}
\begin{center}{\Large\bfseries Score Advanced --- Test 16 (20 May 2023)}\end{center}
\begin{multicols}{2}
\raggedcolumns
\papersec{1}
\partsec{1}{PHYSICS}
\sect{SECTION-I (i)}
\qn{1} A long cylindrical wire of radius `a' carries a current $i$ distributed uniformly over its cross-section. If the magnetic field at distance $r$ and $R$ from the axis have equal magnitude, then :-
\begin{qoptions}
\opt{A}{$a=\dfrac{R+r}{2}$}
\opt{B}{$a=\sqrt{Rr}$}
\opt{C}{$a=Rr/(R+r)$}
\opt{D}{$a=R^2/r$}
\end{qoptions}
\qn{2} A uniform rope of mass $M$ and length $L$ is pulled upwards along a frictionless inclined plane of angle $\theta$ (by applying a force $F(>Mg)$) parallel to incline. The rope is lying along the line of maximum slope on the incline. Tension in the chain at a distance $x$ from the end at which force is applied is :-
\begin{qoptions}
\opt{A}{$F\dfrac{x}{L}+\dfrac{x}{L}Mg\sin\theta$}
\opt{B}{$F\left(1-\dfrac{x}{L}\right)$}
\opt{C}{$F\left(1-\dfrac{x}{L}\right)+\dfrac{x}{L}Mg\sin\theta$}
\opt{D}{None of these}
\end{qoptions}
\qn{3} A uniform disc of mass $m$ has surface distributed charge varying with distance $r$ from the centre as $\sigma(r)=\dfrac{\sigma_0 r}{R}$. It is rotating with angular speed $\omega$ about an axis normal to disc and passing through centre as shown. The ratio of magnitudes of magnetic moment to angular momentum is :-\par
\qfig{p1_phy_s1_q03_1}
\begin{qoptions}
\opt{A}{$\dfrac{\pi\sigma_0R^2}{3m}$}
\opt{B}{$\dfrac{\pi\sigma_0R^2}{2m}$}
\opt{C}{$\dfrac{2\pi\sigma_0R^2}{5m}$}
\opt{D}{$\dfrac{\pi\sigma_0R^2}{4m}$}
\end{qoptions}
\qn{4} Rays of light emitted by a point source S are incident on a uniform ball of transparent material with a refractive index $n$. The ray leaving the source S at an angle $\alpha$ to the straight line MN, on which the source and the center of the ball lie, after two refractions at the boundary of the ball, intersects MN at an angle $\beta$ at point K (see figure). The distance SK is equal to $\ell$. The value of radius $R$ of the ball is given by :- (Take : $\alpha=53^\circ$, $\beta=37^\circ$ and $n=\sqrt2$)\par
\qfig{p1_phy_s1_q04_B}
\begin{qoptions}
\opt{A}{$\dfrac{3\sqrt2\,\ell}{35}$}
\opt{B}{$\dfrac{12\ell}{35}$}
\opt{C}{$\dfrac{6\sqrt2\,\ell}{175}$}
\opt{D}{$\dfrac{\ell}{5}$}
\end{qoptions}
\sect{SECTION-I (ii)}
\qn{5} A sinusoidal wave is propagating along a stretched string that lies along x-axis. The wave is moving in $+x$-direction. Figure shows the graph of transverse displacement of particles of the string at $x=0$ and $x=4$ cm as a function of time. If the particles at $x=0$ and $x=4$ cm are within one wavelength of each other, which of the following is/are correct?\par
\qfig{p1_phy_s2_q05_1}
\begin{qoptions}
\opt{A}{Phase difference between the particles at $x=0$ and $x=4$ cm is $5\pi/6$.}
\opt{B}{Phase difference between the particles at $x=0$ and $x=4$ cm is $7\pi/6$.}
\opt{C}{The speed of wave is 12/5 cm/sec.}
\opt{D}{The speed of wave is 12/7 cm/sec.}
\end{qoptions}
\qn{6} Consider a radioactive decay series $\mathrm{X}\xrightarrow{\alpha\text{-decay}}\mathrm{Y}\xrightarrow{\beta^+\text{-decay}}\mathrm{Z}\xrightarrow{\beta^+\text{-decay}}\mathrm{W}$
\begin{qoptions}
\opt{A}{Y, Z are isobars.}
\opt{B}{X, W are isotones.}
\opt{C}{X, W are isotopes.}
\opt{D}{Z, W are isobars.}
\end{qoptions}
\qn{7} A rod of length $\ell$, very small area of cross section $A$, Young's modulus of elasticity $Y$ is acted upon by two parallel forces $3F$ and $F$ respectively (as shown) and placed on a smooth horizontal plane. If the elastic limit is not crossed and then to study the change in length of rod ($\Delta\ell$) and it's elastic potential energy ($U$) the rod is segmented into four equal parts where magnitude of change in lengths are $\Delta\ell_1$, $\Delta\ell_2$, $\Delta\ell_3$, $\Delta\ell_4$ and elastic potential energy stored in each segment are $U_1$, $U_2$, $U_3$, $U_4$ respectively as shown then which is/are correct ?\par
\qfig{p1_phy_s2_q07_1}
\begin{qoptions}
\opt{A}{$\Delta\ell_1=\dfrac{\Delta\ell_3}{3}=\dfrac{\Delta\ell}{8}$}
\opt{B}{$U_1<\dfrac{U}{4}<U_4$}
\opt{C}{$\Delta\ell_2=\dfrac{\Delta\ell_4}{5}=\dfrac{\Delta\ell}{8}$}
\opt{D}{$U_2<U_3<U_4$}
\end{qoptions}
\qn{8} A uniform circular ring of mass `m' and radius `R' is hinged at its top point `O' and suspended in a vertical plane so that it can freely rotate about a horizontal axis passing through the hinge `O' as shown in the figure. The bead `A' is rigidly attached to the ring and the bead `B' can slide without friction on the ring. The mass of each bead is `m'. The system is released from rest from the initial position shown in the figure. Then choose the correct option(s), immediately after the system is released from rest, (take g= 10m/s$^2$)\par
\qfig{p1_phy_s2_q08_1}
\begin{qoptions}
\opt{A}{The acceleration of bead `A' is $2\sqrt2$ m/s$^2$.}
\opt{B}{The acceleration of bead `A' is $2\sqrt{13}$ m/s$^2$.}
\opt{C}{The acceleration of bead `B' is $2\sqrt2$ m/s$^2$.}
\opt{D}{The acceleration of bead `B' is $2\sqrt{26}$ m/s$^2$.}
\end{qoptions}
\qn{9} Gerasim took a thick rectangular board and drove nails a. b. c. d. e into it. Nails a. b. c are driven into the midpoints of the sides, nails d. e into the corners of the plank. Gerasim attached six pieces of wire to the nails, as shown in the figure. Then he began to slowly and evenly lower the plank into the conducting liquid, applying a potential difference $U$ between the nail a and the liquid. The plank is always located in a vertical plane, keeping de horizontal. The scheme was completely immersed in the liquid during the time $T$. The resistance of each piece of wire is indicated in the figure (the value of $R$ is known). The wire is not covered with insulation; each of its pieces has a constant thickness. The resistance of liquid and nails is negligible. The plate does not conduct electricity.\par
\qfig{p1_phy_s2_q09_1}
\begin{qoptions}
\opt{A}{The current in the circuit remains constant.}
\opt{B}{The current in the circuit for $0<T<\dfrac T2$ is given by $i=\dfrac{U\left(5-\frac{6t}{T}\right)}{R\left(2-\frac{2t}{T}\right)\left(3-\frac{4t}{T}\right)}$.}
\opt{C}{The current in the circuit during $\dfrac T2<t<T$ is given by $i=\dfrac{2U}{R}\left(\dfrac{1}{1-\frac{t-T}{T}}\right)$.}
\opt{D}{The current in the circuit during $\dfrac T2<t<T$ is given by $i=\dfrac{U}{2R}\left(\dfrac{1}{1-\frac{2t-T}{T}}\right)$.}
\end{qoptions}
\qn{10} An electric circuit shown consists of two ideal inductors of inductance $L_1$ and $L_2$, a resistance $R$, a two way switch (S) ($1\rightarrow2$, $1\rightarrow3$), a battery of electromotive force $\varepsilon$ and internal resistance $r$. The switch was in position 2 for a long time and at the instant of time $t=0$, it is suddenly thrown to position 3. Now select correct statement(s) :\par
\qfig{p1_phy_s2_q10_1}
\begin{qoptions}
\opt{A}{Current through the switch `S' as a function of time `t' is $\dfrac{\varepsilon}{r}\left[1-e^{\frac{-Rt}{L_2}}\right]$}
\opt{B}{Current through the switch `S' at any time `t' will be $\dfrac{\varepsilon}{r}$}
\opt{C}{Total heat dissipated in the resistor of resistance $R$ after switch `S' thrown to the position `3' is $\dfrac12\dfrac{L_2\varepsilon^2}{r^2}$}
\opt{D}{Current through resistor of resistance `R' as a function of time is `t' is $\dfrac{\varepsilon}{r}e^{-\frac{Rt}{L_2}}$}
\end{qoptions}
\sect{SECTION-II (i)}
\paraq{1}{2}\textbf{Question Stem}\par
An experiment is performed on thin concavo-convex lens in order to determine the radii $R_1$ (convex) and $R_2$ (concave) of curvature of its two surfaces and the refractive index `$\mu$' of its glass.\par
\textbf{Experiment I:} A real object is placed 36 cm to the left of lens and a real image is formed 72 cm to the right of lens.\par
\textbf{Experiment II:} The lens is floated on the surface of mercury with surface having radius of curvature $R_2$ facing air. A real object is placed at 60 cm above it, and a real image is formed at 30 cm above it (assume mercury to be perfectly reflecting).\par
\textbf{Experiment III:} The lens is floated on the surface of mercury with surface having radius of curvature $R_1$ facing air. A real object is placed at 7.5 cm above it and real image is formed at $\dfrac{20}{3}$ cm above it. On the bases of Experiments I, II and III, answer the following questions.\par
\qn{1} Find the refractive index of material of lens.
\qn{2} Find the magnitude of $R_1$ (in cm)
\paraq{3}{4}\textbf{Question Stem}\par
A circuit shown in the figure in which $k_1$ is closed and $k_2$ is open. Inductor $L$ can be connected to capacitor $C_1$ by closing switch $k_2$ and opening $k_1$ .\par
\qfig{p1_phy_s3_para3_4_1}
\qn{3} Let switch $k_1$ is closed and $k_2$ is opened for long time. Find the charge on capacitor $C_2$ (in $\mu$C)
\qn{4} At $t=0$, when capacitors are fully charged, switch $k_1$ is opened and switch $k_2$ is closed, so that inductor is connected in series with capacitor $C_1$. Current in inductor becomes maximum (for first time) at time $\alpha\pi/2000$. Find the value of $\alpha$.
\paraq{5}{6}\textbf{Question Stem}\par
A particle is projected from a point P (0, 0, 600 m) lying in horizontal X - Z plane with initial velocity of $\vec u=\left(80\hat i+100\hat j-60\hat k\right)$ m/s at $t=0$. It hits a smooth vertical wall in X - Y plane with y-axis in vertically upward direction. Co-efficient of restitution $e=\dfrac12$. (g = 10 m/s$^2$)\par
\qn{5} The ball hits the wall at time $t$ (in sec) when its velocity is making an angle $\theta$ (in degrees) with the vertical. Then write $\theta/t$
\qn{6} Position of the point of impact with the horizontal X-Z plane after collision with the wall is (x,y,z). Find (x + y + z).
\sect{SECTION-II (ii)}
\qn{7} The cam is a circular disk rotating on a shaft that does not pass through the center of the disk. In the manufacture of the cam, a uniform solid cylinder of radius $R$ is first machined. Then an off-center hole of radius $R/2$ is drilled, parallel to the axis of the cylinder, and centered at a point a distance $R/2$ from the center of the cylinder. The cam, of mass $M$, is then slipped onto the circular shaft and welded into place. What is the kinetic energy (in Joule) of the cam when it is rotating with angular speed $\omega$ about the axis of the shaft ? ($M=12$kg, $R=20$ cm, $\omega=\sqrt{200}$ rad/s)\par
\qfig{p1_phy_s4_q07_1}
\qn{8} A sealed balloon, filled with air, floats in water with $\dfrac13$ of its volume submerged. It was found that if it is pushed inside water at a depth $h$, it remains in equilibrium, neither sinking nor rising. Find $h$ (in metres). Given that height of water barometer is 10 m and temperature is constant at all depth.\par
\qfig{p1_phy_s4_q08_1}
\qn{9} In the system shown in the figure, the marking on the bar is equidistant and the blocks $M$ and $m$ are placed on the bar such that the centre of mass of the blocks lies above a mark on the bar. The range of $M$ for which it is possible to balance the weights on the bar is $M_{\min}<M<M_{\max}$. Threads and blocks are weightless. There is no friction. The mass $m$ is known. Value of $\left(\dfrac{M_{\max}+M_{\min}}{2m}\right)$ is :-\par
\qfig{p1_phy_s4_q09_1}
\partsec{2}{CHEMISTRY}
\sect{SECTION-I (i)}
\qn{1} 20 V H$_2$O$_2$ solution is kept in the open vessel and is allowed to decompose by first order kinetics. After 6 hours, 11.35 ml of this H$_2$O$_2$ sample was taken and diluted to 100 ml by adding water. 10 ml of this diluted solution requires 25 ml of 0.025 M KMnO$_4$ in acidic medium for titration. Calculate the rate constant for the decomposition of H$_2$O$_2$. [Given : ln 2 = 0.70, ln 5 = 1.60]
\begin{qoptions}
\opt{A}{0.04 hr$^{-1}$}
\opt{B}{0.05 hr$^{-1}$}
\opt{C}{0.15 hr$^{-1}$}
\opt{D}{0.30 hr$^{-1}$}
\end{qoptions}
\qn{2} At 298 K, $E^\circ_{OP}$ of the following reaction is\par
\centerline{Ag(s) + 2NH$_3$(aq) $\longrightarrow$ Ag(NH$_3$)$_2^+$(aq) + e$^-$}
Given : K$_f$(Ag(NH$_3$)$_2^+$) = 10$^{16}$ ; $E^\circ_{Ag^+|Ag}=0.80$V
\begin{qoptions}
\opt{A}{0.144 V}
\opt{B}{1.744 V}
\opt{C}{$-0.144$ V}
\opt{D}{$-1.744$ V}
\end{qoptions}
\qn{3} The crystal field stabilisation energy (CFSE) for [Co(L)$_6$]$^{2+}$ is 7200 cm$^{-1}$, where L is a monodentate neutral ligand.\par
According to CFT, e$_g$ set of orbitals of central atom of this complex is exactly half filled. Select the correct option regarding this complex :
\begin{qoptions}
\opt{A}{The pairing energy for Co$^{2+}$ is 6000 cm$^{-1}$.}
\opt{B}{The crystal field splitting energy is 4000 cm$^{-1}$.}
\opt{C}{This complex is a high spin complex.}
\opt{D}{This complex absorbs wavelength shorter than that absorbed by [Co(L)$_6$]$^{3+}$.}
\end{qoptions}
\qn{4} Which triene has the largest heat of hydrogenation (magnitude) ?
\begin{qoptions}
\opt{A}{\qfigi{p1_chem_s1_q04_A}}
\opt{B}{\qfigi{p1_chem_s1_q04_B}}
\opt{C}{\qfigi{p1_chem_s1_q04_C}}
\opt{D}{\qfigi{p1_chem_s1_q04_D}}
\end{qoptions}
\sect{SECTION-I (ii)}
\qn{5} Which is/are incorrect statement about refining of copper ?\par
Given impure copper contains impurity of Fe, Ag and Au.\par
$E^\circ_{Ag^+|Ag}=0.799$ V; $E^\circ_{Fe^{+2}|Fe}=-0.44$V; $E^\circ_{Au^{+3}|Au}=1.50$V; $E^\circ_{Cu^{+2}|Cu}=0.34$V
\begin{qoptions}
\opt{A}{Ag and Au remains in anode mud.}
\opt{B}{Fe, Ag, Au remains in anode mud.}
\opt{C}{Blue colour intensity of CuSO$_4$ remains constant during electrolysis.}
\opt{D}{O$_2$(g) is evolved at anode.}
\end{qoptions}
\qn{6} In the following list of four molecules :\par
(I) CO$_2$\quad (II) NO$_2$\quad (III) SO$_2$\quad (IV) SiO$_2$\par
Identify the correct statement(s) :
\begin{qoptions}
\opt{A}{Only I has a linear shape.}
\opt{B}{Only II are gases at room temperature (25\textdegree C)}
\opt{C}{Only III \& IV are soluble in water.}
\opt{D}{All four are acidic oxides.}
\end{qoptions}
\qn{7} Select the correct statement(s) regarding boron compounds :
\begin{qoptions}
\opt{A}{The total number of sp$^3$ hybrid atoms in borax are greater than 10.}
\opt{B}{The order of B--N bond length follows :\par Inorganic Graphite $>$ Borazine $>$ Borazon}
\opt{C}{In solid orthoboric acid, 6 H-bonds are associated per molecule.}
\opt{D}{Maximum 6-atoms are present in a plane in diborane.}
\end{qoptions}
\qn{8} For which of the following monomers is anionic polymerization a favourable mechanism to form polymer ?
\begin{qoptions}
\opt{A}{CH$_2$ = CH $-$ COOCH$_3$}
\opt{B}{CH$_2$ = CH $-$ CN}
\opt{C}{CH$_2$ = CH $-$ OCOCH$_3$}
\opt{D}{CH$_2$ = C(CH$_3$)$_2$}
\end{qoptions}
\qn{9} Which of the following reaction(s) is/are matched CORRECTLY ?
\begin{qoptions}
\opt{A}{\qfigi{p1_chem_s2_q09_A}}
\opt{B}{\qfigi{p1_chem_s2_q09_B}}
\opt{C}{\qfigi{p1_chem_s2_q09_C}}
\opt{D}{\qfigi{p1_chem_s2_q09_D}}
\end{qoptions}
\qn{10} The boiling point of water in a 0.1 molal aqueous barium chloride solution (solution \textbf{A}) is $\mathbf{x}$ \textdegree C. To this solution \textbf{A}, an equal volume of 0.2 molal aqueous silver nitrate solution is added to make a new solution \textbf{B}. The difference in the boiling points of water in the two solutions \textbf{A} and \textbf{B} is $\mathbf{y}\times10^{-2}$ \textdegree C. (Assume : Densities of the solutions \textbf{A} and \textbf{B} are the same as that of water and the soluble salts dissociate completely.)\par
Use: Molal elevation constant (Ebullioscopic Constant), K$_b$ = 0.5 K kg mol$^{-1}$; Boiling point of pure water as 100\textdegree C.)\par
Correct option(s) is/are :
\begin{qoptions}
\opt{A}{$x=100.15$}
\opt{B}{$x=100.10$}
\opt{C}{$|y|=7.5$}
\opt{D}{$|y|=0.15$}
\end{qoptions}
\sect{SECTION-II (i)}
\paraq{1}{2}100 ml 0.1 M weak acid HA (K$_a$ = 10$^{-6}$) is titrated with V ml 0.1 M NaOH (aq).\par
[Given : log 3 = 0.48, log 2 = 0.30]\par
\qn{1} The pH of final solution when V = 25 ml is
\qn{2} The pH of final solution when V = 100 ml is
\paraq{3}{4}\par
X + O$_2$ $\xrightarrow{\text{Catalyst}}$ Y(neutral oxide) + Z(neutral oxide)\par
Z $\xrightarrow{\text{electrolysis}}$ A(inflammable gas) + B(supporter of combustion)\par
Y + B $\longrightarrow$ C(gas)\par
C(gas) + Z $\longrightarrow$ E + D\par
E $\xrightarrow{\text{disproportionation}}$ D + Y + B\par
\qn{3} The total number of lone pair present in X, Y$^+$ and C is-
\qn{4} Let sum of molecular weights of Z, Y and D be M. Find the value of $\dfrac{M-1}{10}$.
\paraq{5}{6}Answer the following questions based on IUPAC nomenclature\par
\qn{5} \qfig{p1_chem_s3_q05_1}
If the IUPAC name of the above compound is $x$-Methyloct-$y$-en-$z$-yne, calculate the value of $x/(y+z)$
\qn{6} \qfig{p1_chem_s3_q05_2}
If the IUPAC name of the above compound is $x,y$-Dimethylcyclohex-$z$-ene, calculate the value of $(x+z)/y$.
\sect{SECTION-II (ii)}
\qn{7} Calculate $\Delta_rG$ (in kJ/mol) for X ($l$, 1 atm, 200 K) $\longrightarrow$ X(g, 1 atm, 200 K)\par
Given : $\Delta_{vap}H$ ($l$, 400 K) = 100 KJ mol$^{-1}$ ; C$_{pm}$(X, $l$) = 180 JK$^{-1}$ mol$^{-1}$\par
C$_{pm}$ (X,g) = 80 KJ mol$^{-1}$ ; Boiling point of X($l$) = 400 K, ln 2 = 0.7
\qn{8} Find the total number of trichlorides from given below which produces tribasic oxy-acid as final product on hydrolysis :\par
BCl$_3$, NCl$_3$, PCl$_3$, AsCl$_3$, SbCl$_3$, BiCl$_3$, FeCl$_3$
\qfullbegin
\qn{9} In how many of the following compounds is the rate of nitration less than the rate of nitration of benzene.\par
\qfig{p1_chem_s4_q09_1}
\qfullend
\partsec{3}{MATHEMATICS}
\sect{SECTION-I (i)}
\qn{1} Let $A$ and $B$ are non-null square matrices of order 2 with real elements, such that $AB=A^2B^2-(AB)^2$ and $\det(B)=2$. Then the value of $\det(A+2B)-\det(B+2A)$
\begin{qoptions}
\opt{A}{4}
\opt{B}{6}
\opt{C}{0}
\opt{D}{1}
\end{qoptions}
\qn{2} For $x\in[0,\pi/2]$, $y=\dfrac{\sin^4x-\cos^4x+\sin^2x\cos^2x}{\sin^4x+\cos^4x+\sin^2x\cos^2x}$, then $y\in$
\begin{qoptions}
\opt{A}{[1/2, 1]}
\opt{B}{[$-$3 /2, 1/ 2]}
\opt{C}{[$-$1, 1]}
\opt{D}{[0,1/2]}
\end{qoptions}
\qn{3} A parabola with latus rectum `4a' slides such that it touches the positive co-ordinate axes. The locus of its focus is
\begin{qoptions}
\opt{A}{$\dfrac{1}{x^2}+\dfrac{1}{y^2}=\dfrac{1}{a^2},\ x>0,y>0$}
\opt{B}{$\dfrac{1}{x^2}+\dfrac{1}{y^2}=\dfrac{2}{a^2},\ x>0,y>0$}
\opt{C}{$\dfrac{1}{x^2}+\dfrac{1}{y^2}=\dfrac{1}{2a^2},\ x>0,y>0$}
\opt{D}{$\dfrac{1}{x^2}+\dfrac{1}{y^2}=\dfrac{4}{a^2},\ x>0,y>0$}
\end{qoptions}
\qn{4} If\par\resizebox{\linewidth}{!}{$\left(2\sin^2x-2\sqrt3\sin x\cos x+1\right)\left(9-4\cos^2x-4\cos x\right)=40$}\par then number of values of $x$ in $[0,20]$ is
\begin{qoptions}
\opt{A}{3}
\opt{B}{4}
\opt{C}{5}
\opt{D}{6}
\end{qoptions}
\sect{SECTION-I (ii)}
\qn{5} Let there be a continuous function $f:[0,1]\rightarrow R$ satisfying $\displaystyle\int_0^1f(x)\,dx=\dfrac13+\int_0^1f^2(x^2)\,dx$. Then
\begin{qoptions}
\opt{A}{$f(x)<x\ \forall\ x\in[0,1]$}
\opt{B}{There is a unique function satisfying this equation}
\opt{C}{$\displaystyle\int_0^1xf(x)\,dx>\int_0^{1/2}f(x)\,dx$}
\opt{D}{$f(1)+f(4)=3$}
\end{qoptions}
\qn{6} For non-negative integers $a$ and $b$, let\par
\resizebox{\linewidth}{!}{$U(a,b)=\dbinom6a\dbinom6b\dbinom{6}{a+b}$}\par
\resizebox{\linewidth}{!}{$T(a,b)=\dbinom6a\dbinom6b\dbinom{6}{|a-b|}$}\par
where $\dbinom nr=\begin{cases}\dfrac{n!}{r!(n-r)!},&r\leqslant n\\ 0,&r>n\end{cases}$ $r$ and $n$ are non-negative integers.\par
Let $S$ denote the sum of all $U(a,b)$.\par
Which of the following statement(s) is(are) CORRECT ?
\begin{qoptions}
\opt{A}{Remainder when $S$ is divided by 1000 is 564}
\opt{B}{Remainder when $S$ is divided by 100 is 84}
\opt{C}{$T(a,3)$ is minimum for two distinct values of $a$}
\opt{D}{$T(a,3)$ is maximum for four distinct values of $a$}
\end{qoptions}
\qn{7} Which of the following is/are \textbf{TRUE}
\begin{qoptions}
\opt{A}{Let $f(x)=\left(x^2+3x+2\right)^{\cos\pi x}$, then the sum of all positive integers `n' for which $\left|\displaystyle\sum_{k=1}^n\log_{10}f(k)\right|=1$ is 54}
\opt{B}{The smallest positive integer with exactly 24 divisors is 360}
\opt{C}{Number of solution(s) of equation $\log_{[x]}(5-[x])=|x-4|$ is 3 (where [.] is GIF)}
\opt{D}{$\sec11^\circ\cdot\sec19^\circ-2\cot71^\circ=2\tan11^\circ$}
\end{qoptions}
\qn{8} Consider the region $R:x^2\leq y^2\leq2-x$, then which of the following is/are \textbf{FALSE} ?
\begin{qoptions}
\opt{A}{Maximum value of $x^2+y^2$ is 8, $\forall(x,y)\in R$}
\opt{B}{Number of integral points (both coordinates integer) in $R$ is 9}
\opt{C}{Area of $R$ is $\dfrac{13}{6}$ sq.units}
\opt{D}{Minimum value of $x^2+y^2$ is 4, $\forall(x,y)\in R$}
\end{qoptions}
\qn{9} Equation of a straight line on the complex plane passing through a point P denoting the complex number $\alpha$ and perpendicular to the line segment OP where `O' is the origin can be written as
\begin{qoptions}
\opt{A}{$\operatorname{Im}\left(\dfrac{z-\alpha}{\alpha}\right)=0$}
\opt{B}{$\operatorname{Re}\left(\dfrac{z-\alpha}{\alpha}\right)=0$}
\opt{C}{$\operatorname{Re}(\bar\alpha z)=0$}
\opt{D}{$\bar\alpha z+\alpha\bar z-2|\alpha|^2=0$}
\end{qoptions}
\qn{10} Consider a series of `n' concentric circles $C_1,C_2,C_3\ldots C_n$ with radii $r_1,r_2\ldots r_n$ respectively, such that $r_1>r_2>r_3>\ldots..>r_n$ and $r_1=20$. If the tangents drawn from any point on $C_i$ to $C_{i+1}$ are such that the chord of contact is a tangent to $C_{i+2}$ ($i=1,2,3\ldots.$) and the angle between the tangents from any point on $C_1$ to $C_2$ is $\dfrac\pi3$ then
\begin{qoptions}
\opt{A}{$r_1,r_2,r_3\ldots\ldots r_n$ are in G.P with common ratio $\dfrac{\sqrt3}{2}$}
\opt{B}{$\displaystyle\lim_{n\to\infty}\sum_{i=1}^nr_i=80+40\sqrt3$}
\opt{C}{$\displaystyle\lim_{n\to\infty}\sum_{i=1}^nr_i=40$}
\opt{D}{$r_2=10\sqrt3$}
\end{qoptions}
\sect{SECTION-II (i)}
\paraq{1}{2}Given two curves $y=f(x)$ passing through the point (0, 1) and $y=g(x)=\displaystyle\int_{-\infty}^xf(t)dt$ passing through the point $\left(0,\dfrac12\right)$. The tangents drawn to both the curves at the points $(x_0,f(x_0))$ and $(x_0,g(x_0))$ intersect on the x-axis for all $x_0$ in the domain, then\par
\qn{1} The value of $f\left(\dfrac12\right)$ is
\qn{2} $\displaystyle\lim_{x\to0}\dfrac{x}{(f(x))^2-1}=$
\paraq{3}{4}Let $X=\{1,2,3\}$. A number is selected from the set $X$ with replacement and this process is repeated 1100 times. For each $i$, $i\in\{1,2,3\}$, let $f(i)$ represent the number of times the number `i' is selected. Also, let $S$ denote the total sum of 1100 numbers selected. If $S^3=162\,f(1)\cdot f(2)\cdot f(3)$, answer the following questions.\par
\qn{3} The harmonic mean of $f(1),f(2),f(3)$ is equal to
\qn{4} There are $f(1)$ red balls, $f(2)$ blue balls, $f(3)$ white balls and `x' green balls (balls of the same color are identical). The balls are arranged in a row such that no two balls of the same color are consecutive. Let $x_1$ and $x_2$ be the minimum and maximum values of $x$ respectively for which the above arrangement is possible. The value of $x_1+x_2$ is equal to
\paraq{5}{6}A person repeatedly tosses a fair dice. He gets two points for a throw of a perfect square and one point otherwise. Let $P_n$ denote the probability that he reaches a score of exactly $n$.\par
\qn{5} Find the value of $\dfrac{P_{10}-P_9}{P_8-P_9}$
\qn{6} Find the value of $\displaystyle\lim_{n\to\infty}P_n$
\sect{SECTION-II (ii)}
\qn{7} The minimum value of $|\sin x+\cos x+\tan x+\sec x+\cot x+\operatorname{cosec}x|$ is $\sqrt k-1$ , then find $k$.
\qn{8} If $y(x^2+3x+2)=1$, then the value of $d^8y/dx^8$ at $x=0$ is $\lambda$. Find $[\lambda\times2^9/8!\,]$, (where [.] denotes greatest integer).
\qn{9} If ABCDEFGH is a parallelopiped and $M$ and $M'$ are the point of intersections of diagonals of parallelograms ABCD and EFGH. If $O$ is a point on $MM'$ such that $\overrightarrow{OA}+\overrightarrow{OB}+\overrightarrow{OC}+\overrightarrow{OD}=\overrightarrow{M'O}$ and $\lambda\overrightarrow{MO}=\overrightarrow{OM'}$, then $\lambda$ is \underline{\hspace{1.5cm}}.
\raggedcolumns
\papersec{2}
\partsec{1}{PHYSICS}
\sect{SECTION-I (i)}
\qn{1} In the arrangement shown in figure pulleys and strings are ideal. End A of string connected to pulley P$_1$ is moved upwards with acceleration $a_A=2$ m/s$^2$ while end B of another string shown in figure is moved up with acceleration $a_B=1$ m/s$^2$. Block C of mass 1 kg is moving up with acceleration 1 m/s$^2$. If block D to which strings are connected symmetrically moves such that its orientations remains same then :- (Assume g = 10 m/s$^2$)\par
\qfig{p2_phy_s1_q01_1}
\begin{qoptions}
\opt{A}{acceleration of block D is $\dfrac12$m/s$^2$ upwards.}
\opt{B}{end A is pulled with force of 22 N.}
\opt{C}{mass of block B is $\dfrac12$ kg.}
\opt{D}{acceleration of block B is $\dfrac73$ m/s$^2$ upwards.}
\end{qoptions}
\qn{2} Select the statement(s) among the following which are \textbf{correct} according to kinetic theory of gas? [All symbols have their usual meaning]
\begin{qoptions}
\opt{A}{A molecule can have a speed greater than $\sqrt{\dfrac{6RT}{M}}$.}
\opt{B}{When the temperature of a gas is doubled, then the most probable speed is twice.}
\opt{C}{Kinetic energy of one mole of monoatomic gas is $\dfrac{3RT}{2}$.}
\opt{D}{Kinetic energy of two moles of diatomic gas is $\dfrac{5RT}{2}$.}
\end{qoptions}
\qn{3} A strong bar magnet of mass $M$ is dropped on a block of copper ($\mu_{rel}\approx1$) from a height of $H$. North pole of the magnet is always upward. Select the correct statement(s) among the following.
\begin{qoptions}
\opt{A}{Initial acceleration of magnet is $\vec g$.}
\opt{B}{Speed of the magnet when it reaches very close to the block of copper is approximately equal to zero.}
\opt{C}{Speed of the magnet when it reaches very close to the block of copper is nearly $\sqrt{2gH}$.}
\opt{D}{The copper block gets heated by nearly $MgH$.}
\end{qoptions}
\qn{4} A small object O is placed symmetrically between two small concave mirrors of same radius of curvature $R$. Find the value(s) of the distance $d$ between the mirrors for which an image formed between the mirrors coincides with O.\par
\qfig{p2_phy_s1_q04_1}
\begin{qoptions}
\opt{A}{$R$}
\opt{B}{1.5$R$}
\opt{C}{2$R$}
\opt{D}{3$R$}
\end{qoptions}
\qn{5} The variation of the current $i$ with time is as shown in the graph. Over one time period,\par
\qfig{p2_phy_s1_q05_1}
\begin{qoptions}
\opt{A}{the r.m.s. value of current is $\sqrt{\dfrac{i_0^2+i_1^2}{2}}$.}
\opt{B}{the r.m.s. value of current is $\sqrt{\dfrac{2i_0^2+i_1^2}{2}}$.}
\opt{C}{average value of current is $\dfrac{i_0+i_1}{2}$.}
\opt{D}{average value of current is $i_0+\dfrac{2i_1}{\pi}$.}
\end{qoptions}
\qn{6} A ring of mass $M$, radius $R$, cross section area $A$ and Young's modulus $Y$ is kept on a smooth cone of base radius $2R$ and semi vertex angle $45^\circ$, as shown in figure. Assume that the extension in the ring is small. Pick the \textbf{correct} statement(s) among the following at equilibrium.\par
\qfig{p2_phy_s1_q06_1}
\begin{qoptions}
\opt{A}{Tension in the ring is $Mg$.}
\opt{B}{Tension in the ring is $\dfrac{Mg}{2\pi}$.}
\opt{C}{Elastic energy stored in the ring is $\dfrac{M^2g^2R}{4\pi YA}$.}
\opt{D}{Loss of gravitational potential energy of ring is $\dfrac{M^2g^2R}{4\pi YA}$.}
\end{qoptions}
\sect{SECTION-I (ii)}
\paraq{7}{8}A soap film is made inside a fixed circular frame (centre at A), a fixed rod AB and a uniform movable rod AC having a mass $M$ and length $R$. Radius of frame is $R$ and surface tension of the liquid is $S$. The rod AC is free to rotate about an axis through point A and perpendicular to the plane of the frame. Initially the ends B and C are very close to each other. Neglect gravity and friction.\par
\qfig{p2_phy_s2_para7_8_1}
\qn{7} Find the maximum value of the angular speed of the rod after it is released. Assume rod is much heavier than the liquid in the film.
\begin{qoptions}
\opt{A}{$\sqrt{\dfrac{6\pi S}{M}}$}
\opt{B}{$\sqrt{\dfrac{12\pi S}{M}}$}
\opt{C}{$\sqrt{\dfrac{24\pi S}{M}}$}
\opt{D}{$\sqrt{\dfrac{36\pi S}{M}}$}
\end{qoptions}
\qn{8} Now assume that the frame and the rod AC are highly conducting and the fixed rod AB has a resistance $r$. A uniform magnetic field of magnitude $B$ directed perpendicular to the frame is present in the region. Now the rod AC is released at same position as in the previous part. It is found that the rod quickly gains angular speed $\omega$ which remains constant throughout the remaining motion till end C reaches B. Then, [Liquid is non-conducting]
\begin{qoptions}
\opt{A}{$\omega=\dfrac{2Sr}{B^2R^2}$}
\opt{B}{$\omega=\dfrac{4Sr}{B^2R^2}$}
\opt{C}{$\omega=\dfrac{8Sr}{B^2R^2}$}
\opt{D}{$\omega=\dfrac{12Sr}{B^2R^2}$}
\end{qoptions}
\paraq{9}{10}A point charge of mass $m$ and charge $q$ is placed as shown near an infinitely large non conducting plane having a circular cavity of radius $R$. The point charge is at a perpendicular distance $R$ from centre of cavity. The plane has a uniform charge density $\sigma$.\par
\qfig{p2_phy_s2_para9_10_1}
\qn{9} Flux of electrostatic field produced by the point charge through the surface of the non-conducting plane is :-
\begin{qoptions}
\opt{A}{$\dfrac{q}{2\varepsilon_0}$}
\opt{B}{$\dfrac{q}{\sqrt2\varepsilon_0}$}
\opt{C}{$\dfrac{q}{2\varepsilon_0}\left(1-\dfrac{1}{\sqrt2}\right)$}
\opt{D}{$\dfrac{q}{2\sqrt2\varepsilon_0}$}
\end{qoptions}
\qn{10} If the particle is released then the initial rate of change of its speed is :-
\begin{qoptions}
\opt{A}{zero}
\opt{B}{$\dfrac{q\sigma}{\sqrt2\varepsilon_0m}$}
\opt{C}{$\dfrac{q\sigma}{2\varepsilon_0m}\left(1-\dfrac{1}{\sqrt2}\right)$}
\opt{D}{$\dfrac{q\sigma}{2\sqrt2\varepsilon_0m}$}
\end{qoptions}
\sect{SECTION-II (i)}
\paraq{1}{2}\textbf{Stem-I}\par
A sample of He$^+$ ions is excited by photons of energy 48.36 $eV$. Photons are emitted when the ions return to their ground state. This radiation falls on a metallic plate of work function $6eV$. $\left[h=6.6\times10^{-34}\text{J.s},\ m_e=9.1\times10^{-31}\text{kg}\right]$\par
\qn{1} The maximum wavelength of photon emitted when the He$^+$ ions return to their ground state is approximately \underline{\hspace{1.5cm}} nm.
\qn{2} The minimum de-Broglie wavelength of the photoelectron emitted is approximately\underline{\hspace{1.5cm}} \AA.
\qfullbegin
\paraq{3}{4}\textbf{Stem-II}\par
\qfig{p2_phy_s3_para3_4_1}
\qfullsplit
You have a microammeter that reads 50mA at full-scale deflection and the coil in the meter movement has a resistance of 20 ohms. By adding two resistors. $R_1$ and $R_2$ and a 1.5 volt battery as shown in figure, you can convert this into an ohmmeter. When the two out-coming leads of this ohmmeter are connected together, the meter is to resister zero ohms by giving exactly full-scale deflection. When the leads are connected across an unknown resistance $R$, the deflection will indicate the resistance value if the scale is appropriately calibrated. If we want a reading of 15 ohm at the half of the scale reading :\par
\qfullend
\qn{3} The approximate value (in $\Omega$) of $R_2$ is :-
\qn{4} The approximate value (in $\Omega$) of $R_1$ is :-
\paraq{5}{6}\textbf{Stem = III}\par
There are two Vernier callipers both of which have 1 $cm$ divided into 10 equal divisions on the main scale. The Vernier scale of one of the callipers C$_1$ has 10 equal divisions that correspond to 9 main scale divisions. The Vernier scale of the other callipers C$_2$ has 10 equal divisions that correspond to 11 main scale divisions. The readings of the two callipers are shown in the figure.\par
\qfig{p2_phy_s3_para5_6_1}
\qn{5} Reading of callipers C$_1$ is \underline{\hspace{1.5cm}} $cm$.
\qn{6} Reading of callipers C$_2$ is \underline{\hspace{1.5cm}} cm.
\sect{SECTION-II (ii)}
\qn{7} In a young's double slit experiment P is the second point from central maximum where Intensity is $\dfrac34$th of maximum intensity. The phase difference between the waves from the two slits reaching at that point is $\dfrac{n\pi}{24}$. Find the value of $n$.
\qn{8} Determine the resistance (in ohm) to be connected across the points A, B in given circuit, to dissipate maximum power through it.\par
\qfig{p2_phy_s4_q08_1}
\qn{9} At the vertices of the regular 20-gon, there are 20 identical lenses. The optical centers of the lenses are located exactly at the vertices of the polygon, the planes of all lenses are perpendicular to one of the sides adjacent to the lens. The focal lengths of the lenses are all equal to $F=10$ cm and coincide with the length of the side of the 20-gon. One of the lenses is illuminated by a parallel light flux directed along its optical axis. It turns out that one of the rays has a closed trajectory. Determine the radius of the circle (in cm) inscribed in this trajectory. Consider all of the lenses are converging and all angles small such that $\sin\alpha\approx\tan\alpha\approx\alpha$.\par
\qfig{p2_phy_s4_q09_1}
\raggedcolumns
\partsec{2}{CHEMISTRY}
\sect{SECTION-I (i)}
\qn{1} Molar conductance of 0.012 M aqueous solution of Cl--CH$_2$--C(=O)--OH acid is 100 $\Omega^{-1}$ cm$^2$ mol$^{-1}$. The ion conductance of Cl--CH$_2$--C(=O)--O$^-$ and H$^+$ ions are 50 $\Omega^{-1}$ cm$^2$ mol$^{-1}$ and 350 $\Omega^{-1}$ cm$^2$ mol$^{-1}$, respectively at infinite dilution. (use log 3 = 0.48)
\begin{qoptions}
\opt{A}{degree of dissociation = 0.25}
\opt{B}{[H$^+$] = 3$\times$10$^{-3}$ M}
\opt{C}{pK$_a$ of Cl--CH$_2$--C(=O)--OH = 3}
\opt{D}{pH = 2.52}
\end{qoptions}
\qn{2} Choose the \textbf{CORRECT} statement(s) regarding Vanilin\par
\qfig{p2_chem_s1_q02_1}
\begin{qoptions}
\opt{A}{\textbf{[P]} can be obtained as a major product when 3,4-dihydroxy benzaldehyde is treated with one equivalent of NaOH and one equivalent of MeBr.}
\opt{B}{\textbf{[P]} can be obtained as a one of the product when 2-methoxy phenol is treated with chloroform and aq KOH.}
\opt{C}{Positional isomer of \textbf{[P]} can be obtained when 3,4-dihydroxy benzaldehyde is treated with one equivalent of NaOH followed by treatment with one equivalent of MeBr.}
\opt{D}{\textbf{[P]} can be obtained when 3,4-dihydroxy benzaldehyde is treated with two equivalent of NaOH and one equivalent of MeBr followed by acidic work up by dilute HCl.}
\end{qoptions}
\qn{3} Analyze the following data and choose the \textbf{CORRECT} option(s)\par
{\footnotesize\setlength{\tabcolsep}{2pt}\noindent\begin{tabular}{|L{3.9cm}|L{3.9cm}|}\hline
\centering\textbf{Compound [X]} & \centering\textbf{Compound [Y]}\tabularnewline\hline
M.F : C$_4$H$_6$O (acyclic compound) & M.F : C$_4$H$_6$O (acyclic compound)\\\hline
Gives positive tollens test & Gives positive tollens test\\\hline
Gives positive bromine water test & Gives positive bromine water test\\\hline
Can show tautomerism & Cannot show tautomerism\\\hline
\end{tabular}}\par
\begin{qoptions}
\opt{A}{\textbf{[X]} \& \textbf{[Y]} are positional isomers}
\opt{B}{\textbf{[Y]} is an aldol condensation product of acetaldehyde}
\opt{C}{\textbf{[X]} on treatment with $\Delta_3$O$^+$, number of replaceable hydrogens with Deuterium are 4}
\opt{D}{\textbf{[X]} \& \textbf{[Y]} are chain isomers}
\end{qoptions}
\qn{4} Which of the following statement/comparisons are incorrect ?
\begin{qoptions}
\opt{A}{K$_4$[FeCl(CN)$_4$(O$_2$)] can be diamagnetic if central atom is d$^2$sp$^3$ hybridised}
\opt{B}{Ni(CO)$_4$ is high spin while [Ni(CN)$_4$]$^{2-}$ is low spin}
\opt{C}{O--O bond length increases in the following oxidation reaction :\par Na$_4\overset{\text{II}}{[\text{Zn}]}$(CN)$_4$(O$_2$)$_2$] $\longrightarrow$ Na$_2\overset{\text{II}}{[\text{Zn}]}$(CN)$_4$(O$_2$)$_2$]}
\opt{D}{CuSO$_4$ + KCN(excess) $\longrightarrow$ X(Cu containing) + Other product. There are 5 atoms lying in same plane in anionic part of X}
\end{qoptions}
\qn{5} Choose the correct statement(s) :
\begin{qoptions}
\opt{A}{Sulphur, phosphorus and chlorine are all sp$^3$ hybridised in H$_2$S$_2$O$_7$, H$_3$PO$_3$ and HClO$_3$ respectively}
\opt{B}{XeF$_4$ + SbF$_5$ $\longrightarrow$ Y (ionic compound)\par The hybridisation of central atom in cationic part of Y is sp$^3$d}
\opt{C}{In ClF$_5$ , the lone pair of central atom is present in axial orbital.}
\opt{D}{Fe + HCl $\longrightarrow$ X + Y (g)\par X is used as halogen carrier in chlorination of benzene}
\end{qoptions}
\qn{6} A compound AB$_2$ crystallizes similarly as CaF$_2$. $r_{A^{+2}}=$ `x' $r_{B^-}=$ `y' where $r_{A^{+2}}$ and $r_{B^-}$ are respective radii. Select the correct statement(s) for the above crystal :
\begin{qoptions}
\opt{A}{Minimum distance between two B$^-$ ions will be $\dfrac{2}{\sqrt3}(x+y)$}
\opt{B}{Two B$^-$ ions may be present at a distance of $2(x+y)$}
\opt{C}{One A$^{+2}$ ion is surrounded by eight B$^-$ ions in lattice}
\opt{D}{It can show anion Frenkel defect}
\end{qoptions}
\sect{SECTION-I (ii)}
\paraq{7}{8}\par
\qfig{p2_chem_s2_para7_8_1}
\qn{7} Choose the \textbf{INCORRECT} statement
\begin{qoptions}
\opt{A}{\textbf{[P]} \& \textbf{[Q]} are obtained from toluene by electrophilic substitution reaction}
\opt{B}{SO$_3$H will get substituted on benzene ring in first step and product formed is paramethylbenzenesulphonic acid and ortho-methylbenzenesulphonic acid.}
\opt{C}{\textbf{[P]} and \textbf{[Q]} are positional isomers}
\opt{D}{SO$_2$Cl will get substituted in first step}
\end{qoptions}
\qn{8} Choose the \textbf{CORRECT} Option
\begin{qoptions}
\opt{A}{\textbf{[P]} is the ortho isomer}
\opt{B}{\textbf{[Q]} is the ortho isomer}
\opt{C}{Number of oxygen atoms in \textbf{[S]} are 3}
\opt{D}{Number of sp$^2$ carbon atoms in saccharine are 6}
\end{qoptions}
\paraq{9}{10}Thermodynamics stability is often used in chemical reactions to predict extent of reactions or to predict the prefferred product in some reaction. Thermodynamically a more stable substance will be the one which has lower Gibb's free energy. In various cases it might be decided on the basis of enthalpy, however that may not always be true. The thermodynamic data therefore can be used for predicting stability among isotopes \& to decide the possible transformation. Based on this information \& the data given below, answer the questions that follow.\par
\textbf{Data:}\quad $\Delta_fH^\circ$ of Sn$_{(\text{white})}=0$\qquad $\Delta_fH^\circ$ of Sn$_{(\text{gray})}=-2.1$ kJ/mole\par
\phantom{\textbf{Data:}}\quad S$^\circ$ of Sn$_{(\text{white})}=52$ J/K-mole\qquad S$^\circ$ of Sn$_{(\text{gray})}=44$ J/K-mole\par
All above data at 300 K.\par
Density of Sn$_{(w)}=7.36$ g/ml\par
Density of Sn$_{(g)}=5.769$ gm/ml\par
One such transformation is `tin pest' which is an autocatalytic, allotropic transformation of the element tin causing deterioration of tin objects at low temperatures.\par
\qn{9} Which of the following graph is correct for the allotropic transition Sn$_{(\text{white})}\rightarrow$ Sn$_{(\text{gray})}$ at 1 bar \& 300 K.
\begin{qoptions}
\opt{A}{\qfigi{p2_chem_s2_q09_A}}
\opt{B}{\qfigi{p2_chem_s2_q09_B}}
\opt{C}{\qfigi{p2_chem_s2_q09_C}}
\opt{D}{\qfigi{p2_chem_s2_q09_D}}
\end{qoptions}
\qn{10} As temperature is reduced, what can be predicted for objects containing tin
\begin{qoptions}
\opt{A}{There will be a expansion in volume observed due to the allotropic change.}
\opt{B}{There will be a contraction in volume observed due to the allotropic change.}
\opt{C}{No change takes places}
\opt{D}{Expansion or compression both may occur due to allotropic change}
\end{qoptions}
\sect{SECTION-II (i)}
\paraq{1}{2}Percentage yield is very critical for any organic reaction either involving subtitution or elimination. If mass of product obtained and reactant initially taken is precisely quoted then the balanced stoichiometric equation can be used to compute percentage yield of reaction.\par
percentage yield = (actual yield / theoretical yield) $\times$ 100\par
\qn{1} If $X$ is molecular weight of the major organic product obtained in the given reaction then what is $X$ / 7 (C = 12.0107 amu, H = 1.0079 amu, O = 15.9994 amu, Br =78.9183 amu)\par
\qfig{p2_chem_s3_q01_1}
\qn{2} Consider the reaction sequence from \textbf{P} to \textbf{Q} shown below, the overall yield of the major product \textbf{Q} from \textbf{P} is 75\%. What is the amount in grams of \textbf{Q} obtained from 4.65 ml of \textbf{P}? (Use density of P = 1 g ml$^{-1}$)\par
\qfig{p2_chem_s3_q02_1}
\paraq{3}{4}A(g) reacts with Cl$_2$(g) in presence of charcoal to form B. P$_4$ reacts with B to form A(g) and C. C reacts with CH$_3$COOH to give D (Organic compound) alongwith other products. 7.85g D reacts with excess ammonia in presence of pyridine to give E (Organic compound)\par
Let weight of E formed = x g\par
Also, 0.1 mole of B on complete hydrolysis gives acids F and G. Let weight of NaOH required to neutralise mixture of F and G be y grams.\par
[Atomic weight : Na = 23, Cl = 35.5, S = 32, P = 31, O = 16, H = 1, C = 12, N = 14, F = 19, Br = 80, I = 127]\par
\qn{3} Find $x$
\qn{4} Find $\left(\dfrac{y}{10}\right)$
\paraq{5}{6}Consider following reagents with given codes :\par
\textbf{Code 1 :} dil.HCl\par
\textbf{Code 2 :} NH$_3$ (aq)\par
\textbf{Code 3 :} KI (aq)\par
\textbf{Code 4 :} KCN (aq)\par
\textbf{Code 5 :} Na$_2$S$_2$O$_3$ (aq)\par
\textbf{Code 6 :} NaOH (aq)\par
Let the sum of codes of reagents in which AgCl is soluble be $x$.\par
Let the sum of codes of reagents in which Zn(OH)$_2$ is soluble be $y$\par
\qn{5} Find $\dfrac{x}{2}$
\qn{6} Find $\dfrac{y}{4}$
\sect{SECTION-II (ii)}
\qn{7} Assume that the reaction of MeMgBr with ethylacetate proceeds with 100\% conversion to give \textit{tert}-butanol. Calculate the volume (in mL) of 0.2 M solution of MeMgBr required to convert 10 mL of a 0.25 M solution of ethylacetate to tert-bulanol.
\qn{8} What is the sum of position number assigned to Bromine atoms while numbering the parent chain in given structures?\par
\qfig{p2_chem_s4_q08_1}
\qn{9} From the given list of sols how many can be used to coagulate the haemoglobin sol by addition to it ?\par
(i) Fe(OH)$_3$ sol\quad (ii) Ca(OH)$_2$ sol\quad (iii) Gold sol\quad (iv) Clay\par
(v) As$_2$S$_3$ sol\quad (vi) Al(OH)$_3$ sol\quad (vii) CdS sol\quad (viii) Cd(OH)$_2$ sol.
\raggedcolumns
\partsec{3}{MATHEMATICS}
\sect{SECTION-I (i)}
\qn{1} Let $A$ and $B$ be two square matrices of order 3 then which of the following statement(s) is(are) correct ?
\begin{qoptions}
\opt{A}{$ABA^T$ is symmetric matrix}
\opt{B}{$AB-BA$ is skew symmetric matrix}
\opt{C}{If $B=|A|A^{-1},|A|\neq0$ then $\operatorname{adj.}(A^T)-B$ is a skew symmetric matrix}
\opt{D}{If $B+A^T=O$ and $A$ is skew symmetric matrix then $B^{15}$ is also skew symmetric matrix}
\end{qoptions}
\qn{2} $L_1$ is a tangent drawn to the curve $x^2-4y^2=16$ at $A\left(5,\dfrac32\right)$. $L_2$ is another tangent parallel to $L_1$ which meets the curve at $B$. $L_3$ and $L_4$ are normals to the curve at $A$ and $B$. Lines $L_1,L_2,L_3,L_4$ forms a rectangle, then
\begin{qoptions}
\opt{A}{Equation of tangent at $B$ is $6y=5x-16$}
\opt{B}{Equation of normal at $B$ is $12x+10y+75=0$}
\opt{C}{Radius of largest circle inscribed in the rectangle is $\dfrac{32}{\sqrt{61}}$}
\opt{D}{Radius of the circle circumscribing the rectangle is $\dfrac{\sqrt{109}}{2}$}
\end{qoptions}
\qn{3} Let $a_1,a_2,a_3\in R-\{0\}$ and $a_1+a_2\cos2x+a_3\sin^2x=0$ for all $x\in R$, then
\begin{qoptions}
\opt{A}{Vectors $\vec a=a_1\hat i+a_2\hat j+a_3\hat k$ and $\vec b=4\hat i+2\hat j+\hat k$ are perpendicular to each other}
\opt{B}{Vectors $\vec a=a_1\hat i+a_2\hat j+a_3\hat k$ and $\vec b=-\hat i+\hat j+2\hat k$ are parallel to each other}
\opt{C}{If vector $\vec a=a_1\hat i+a_2\hat j+a_3\hat k$ is of length $\sqrt6$, then one of the ordered triplet $(a_1,a_2,a_3)=(1,-1,-2)$}
\opt{D}{If $2a_1+3a_2+6a_3=26$, then $\left|a_1\hat i+a_2\hat j+a_3\hat k\right|$ equals $2\sqrt6$}
\end{qoptions}
\qn{4} Let $f:I\rightarrow I$, defined as $f(x)=2\sin(2\pi x)-10\tan(5\pi x)+7\cos(4\pi x)+3$, then which of the following statement(s) is/are \textbf{TRUE} ?\par
[\textbf{Note:} $I$ denote the set of all integers]
\begin{qoptions}
\opt{A}{$f(x)$ is periodic function}
\opt{B}{$f(x)$ is an even function}
\opt{C}{$f(x)$ is an odd function and its inverse exists}
\opt{D}{$f(f(f(x)))=f(f(x))$ for all $x\in I$}
\end{qoptions}
\qn{5} Let $J=\displaystyle\int_0^{1/2}\left(\dfrac14-x^2\right)^4dx$ and $K=\displaystyle\int_0^{1/2}x^4(1-x)^4dx$, then
\begin{qoptions}
\opt{A}{$\dfrac JK=2$}
\opt{B}{$J-K=0$}
\opt{C}{$J=\dfrac12\displaystyle\int_0^1x^4(1-x)^4dx$}
\opt{D}{$K=\dfrac{1}{1260}$}
\end{qoptions}
\qn{6} If a curve passing through P(4, 2) satisfying the differential equation $\dfrac{dy}{dx}=\dfrac{y}{x-y^2}$ is a conic then which of the following statement(s) is(are) correct ?
\begin{qoptions}
\opt{A}{Focus of the conic lies on line $y=2$}
\opt{B}{Equation of directrix of conic is $4x-17=0$}
\opt{C}{Eccentricity of the conic is 1}
\opt{D}{Eccentricity of the conic is 3}
\end{qoptions}
\sect{SECTION-I (ii)}
\paraq{7}{8}Let $y=f(x)$ be a solution of the differential equation $\dfrac{dy}{dx}=-4x(y-2)$ with $f(0)=2$ and $g(x)=\max.\left\{1+\dfrac x2-\left[\dfrac x2\right],\sin x,|x-1|\right\}$, $h(x)=\ln x$\par
[\textbf{Note:} $[y]$ denotes greatest integer function less than or equal to $y$]\par
\qn{7} Area enclosed by the curves $y=f(x)$, $y=g(x)$ and the y-axis, is
\begin{qoptions}
\opt{A}{1}
\opt{B}{2}
\opt{C}{3}
\opt{D}{4}
\end{qoptions}
\qn{8} Area bounded by the curve $y=f(x)$, $y=h(x)$ and the y-axis, is
\begin{qoptions}
\opt{A}{$e$}
\opt{B}{$e^2$}
\opt{C}{$\dfrac{3e}{2}$}
\opt{D}{$e^2-1$}
\end{qoptions}
\paraq{9}{10}Let $f(x)=-x^2+bx+c$ be a function such that $f(x)=f(2-x)$ and $\alpha,\beta$ be the roots of the equation $f(x)=0$.\par
[\textbf{Note:} $[k]$ denotes the greatest integer function less than or equal to $k$]\par
\qn{9} If $c\in(0,1)$ value of $\max([\alpha],[\beta])$ is equal to
\begin{qoptions}
\opt{A}{2}
\opt{B}{4}
\opt{C}{6}
\opt{D}{8}
\end{qoptions}
\qn{10} If $c\in(-1,0)$ then which of the following statement(s) is incorrect ?
\begin{qoptions}
\opt{A}{$\min.([\alpha],[\beta])=0$}
\opt{B}{$\max.([\alpha],[\beta])=2$}
\opt{C}{$[\alpha]+[\beta]<2$}
\opt{D}{Maximum value of $f(x)\in(0,1)$}
\end{qoptions}
\sect{SECTION-II (i)}
\paraq{1}{2}\textbf{Question Stem}\par
Let the line $L\equiv ax+by+c=0$ where $a,b,c\in N$ is tangent to the parabola $y^2=8x$ at a point where product of abscissa and ordinate is equal to $-1$.\par
\qn{1} If the y intercept of line $L$ is equal to $\lambda$ then value of $|\lambda|$ is
\qn{2} The number of tangents to $y=\sin(x+y)$, $x\in[-2\pi,2\pi]$ which are parallel to $L$, is equal to
\paraq{3}{4}\textbf{Question Stem}\par
Let $S_1=\displaystyle\sum_{k=1}^{10}k(k-1)\,{}^{10}C_k$, $S_2=\displaystyle\sum_{k=1}^{10}k\cdot{}^{10}C_k$ and $S_3=\displaystyle\sum_{k=1}^{10}k^2\cdot{}^{10}C_k$\par
\qn{3} The value of $\dfrac{S_1}{S_2}$ is
\qn{4} The value of $\dfrac{S_3}{2^{10}}$ is
\paraq{5}{6}\textbf{Question Stem}\par
Let $f(x)$ and $g(x)$ be two functions satisfying $f(x^4)+g(2-x)=3x^2$ and $g(2-x)+g(x)=0$\par
\qn{5} The value of $\displaystyle\int_1^{-1}g(1+x)\,dx$ is equal to
\qn{6} The value of $\displaystyle\int_{-2}^2f(x^4)\,dx$ is equal to
\sect{SECTION-II (ii)}
\qn{7} The greatest possible integral value of $\dfrac{b-a}{\tan^{-1}b-\tan^{-1}a}$, (where $0<a<b<\sqrt3$ ) is
\qn{8} If $\displaystyle\int_0^{\pi/2}\dfrac{dx}{\left(\sqrt{\cos x}+\sqrt{\sin x}\right)^4}=J$ then $\dfrac1J$ is
\qn{9} If $f(x)=\left(\left[\{x\}\right]\tan^{-1}\left(\dfrac{x^2-3x-1}{x^2-3x+5}\right)+3-x^7\right)^{\frac17}$\par
Where $[k]$ and $\{k\}$ denotes greatest integer and fractional part of $k$ respectively, then the value of $f^{-1}(50)-f(50)+f(f(100))$ , is
\end{multicols}
\end{document}
